\documentclass[11pt]{article}
\usepackage{preamble}
\renewcommand{\answer}[1]{}

\begin{document}

\ladderheading{AP Statistics \textperiodcentered\ Unit 2 \textperiodcentered\ Topic 2.2 \textperiodcentered\ B: 2026-10-16 \textperiodcentered\ E: 2026-10-30}{Which Pass Rate Answers Which Question?}

\focusbanner{TODAY'S PROBLEM}

Suppose a school records exam outcomes for 300 students. Students chose for themselves whether to use the tutoring center.\par\smallskip \textbf{Mara:} \emph{The pass rate is 75\%.}\par \textbf{Noah:} \emph{The pass rate is 88\%.}\par \textbf{Elena:} \emph{Tutoring raises the pass rate by 13 percentage points because $88\%-75\%=13$ percentage points.}\par
\begin{center}
\textbf{Exam outcomes for 300 students}\par\smallskip
\begin{tabular}{|l|c|c|c|}
\hline
\textbf{Center?} & \textbf{Passed} & \textbf{Did not pass} & \textbf{Total} \\ \hline
\textbf{Used} & 66 & 9 & 75 \\ \hline
\textbf{Did not} & 159 & 66 & 225 \\ \hline
\textbf{Total} & 225 & 75 & 300 \\ \hline
\end{tabular}
\end{center}
\begin{firsttakebox}{FIRST TAKE --- before the video (5 min)}
Can 75\% and 88\% both be correct pass rates? State what group you think each percentage describes.
\workspace{0.9in}
\end{firsttakebox}

\begin{callout}{\IconBook\ NAME THE DENOMINATOR, THEN THE CLAIM (keep on the page)}{calloutgreen}
A \textbf{marginal relative frequency} uses a row or column total divided by the table total. A
\textbf{conditional relative frequency} uses a cell count divided by the total for the named group.
To decide whether two categorical variables are associated, compare the same response's conditional
percentage across the groups. Different conditional distributions show association. They do not show cause
when people choose their own groups; a fair causal comparison requires random assignment.
\end{callout}

\newpage
\textbf{Finish after the video:}\par\smallskip

\dokbadge{1}~\textbf{(a)} Label 75 percent and 88 percent as marginal or conditional. Name the group for each.\par
\workspace{0.7in}

\dokbadge{2}~\textbf{(b)} Verify the 75 percent and 88 percent rates. The nonuser pass rate is about 70.7 percent; find the user-minus-nonuser gap.\par
\workspace{1.4in}

\dokbadge{3}~\textbf{(c)} Can Mara and Noah both be right? Explain their denominators. Use the two group rates to correct Elena's 13-point comparison, and explain why her word raises goes beyond these data. Write 3--4 sentences.\par
\workspace{2.0in}

\begin{sentenceframebox}\raggedright\textbf{Frame:} Mara answers \blank[2.1in], while Noah answers \blank[2.1in].\end{sentenceframebox}

\begin{sentenceframebox}\raggedright\textbf{Frame:} The variables appear \blank[0.8in] because \blank[2.1in]; the word \emph{raises} is unwarranted because \blank[2.0in].\end{sentenceframebox}

\begin{summaryexitbox}{EXIT --- turn this in}
One thing the video changed about my first take: \hrulefill\par
\workspace{0.4in}
\end{summaryexitbox}

\end{document}
